% from ega0-1-fr.tex
\oldpage[0\textsubscript{I}]{36}
topologique \emph{sous-jacent} à l'espace annelé $(X,\mathcal{A})$, et
$\mathcal{A}$ le \emph{faisceau structural}. Ce dernier se note
$\mathcal{O}_X$, et sa fibre en un point $x\in X$ se note $\mathcal{O}_{X,x}$
ou simplement $\mathcal{O}_x$ lorsqu'il n'en résulte pas de confusion.

On désignera par $1$ ou $e$ la \emph{section unité} de $\mathcal{O}_X$
au-dessus de $X$ (élément unité de $\Gamma(X,\mathcal{O}_X)$).

Comme dans ce Traité nous aurons surtout à considérer des faisceaux d'anneaux
\emph{commutatifs}, il sera sous-entendu, lorsque nous parlerons d'un espace
annelé $(X,\mathcal{A})$ sans préciser, que $\mathcal{A}$ est un faisceau
d'anneaux commutatifs.

Les espaces annelés à faisceau structural non nécessairement commutatif
(resp. les espaces topologiquement annelés) forment une \emph{catégorie},
lorsqu'on définit un \emph{morphisme}
$(X,\mathcal{A})\to(Y,\mathcal{B})$ comme un couple
$(\psi,\theta)=\Psi$ formé d'une application continue $\psi:X\to Y$ et d'un
$\psi$-\emph{morphisme} $\theta:\mathcal{G}\to\mathcal{F}$ (3.5.1) de
faisceaux d'anneaux (resp. de faisceaux d'anneaux topologiques)~; le
\emph{composé} d'un second morphisme
$\Psi'=(\psi',\theta'):(Y,\mathcal{B})\to(Z,\mathcal{C})$ et de $\Psi$, noté
$\Psi''=\Psi'\circ\Psi$, est le morphisme $(\psi'',\theta'')$ où
$\psi''=\psi'\circ\psi$, et $\theta''$ est le composé de $\theta$ et $\theta'$
(égal à $\psi'_*(\theta)\circ\theta'$, cf. 3.5.2). Pour les espaces
annelés, rappelons que l'on a alors
$\theta''^\sharp=\theta^\sharp\circ\psi^*(\theta'^\sharp)$ (3.5.5)~; donc si
$\theta'^\sharp$ et $\theta^\sharp$ sont des homomorphismes \emph{injectifs}
(resp. \emph{surjectifs}), il en est de même de $\theta''^\sharp$, compte tenu
du fait que $\psi_x\circ\rho_{\psi(x)}$ est un isomorphisme pour tout $x\in X$
(3.7.2). On vérifie aussitôt, grâce à ce qui précède, que lorsque $\psi$
est une application continue \emph{injective} et $\theta^\sharp$ un
homomorphisme \emph{surjectif} de faisceaux d'anneaux, le morphisme
$(\psi,\theta)$ est un \emph{monomorphisme} (T, 1.1) dans la catégorie des
espaces annelés.

Par abus de langage, on remplacera souvent $\psi$ par $\Psi$ dans les
notations, par exemple en écrivant $\Psi^{-1}(U)$ au lieu de $\psi^{-1}(U)$
pour une partie $U$ de $Y$, lorsque cela ne risquera pas d'entraîner confusion.
\end{env}

\begin{env}[4.1.2]
\label{0.4.1.2-fr}
Pour toute partie $M$ de $X$, le couple $(M,\mathcal{A}|M)$ est évidemment un
espace annelé, dit \emph{induit} sur $M$ par l'espace annelé
$(X,\mathcal{A})$ (et appelé encore la \emph{restriction} de
$(X,\mathcal{A})$ à $M$). Si $j$ est l'injection canonique $M\to X$ et
$\omega$ l'application identique de $\mathcal{A}|M$, $(j,\omega^b)$ est un
monomorphisme $(M,\mathcal{A}|M)\to(X,\mathcal{A})$ d'espaces annelés, appelé
l'\emph{injection canonique}. Le composé d'un morphisme
$\Psi:(X,\mathcal{A})\to(Y,\mathcal{B})$ et de cette injection est appelé la
\emph{restriction} de $\Psi$ à $M$.
\end{env}

\begin{env}[4.1.3]
\label{0.4.1.3-fr}
Nous ne reviendrons pas sur la définition des $\mathcal{A}$-Modules ou
faisceaux algébriques sur un espace annelé $(X,\mathcal{A})$
(G, II, 2.2)~; lorsque $\mathcal{A}$ est un faisceau d'anneaux non
nécessairement commutatifs, par $\mathcal{A}$-Module, il faudra toujours
sous-entendre \og $\mathcal{A}$-Module à gauche\fg{} sauf mention expresse du
contraire. Les sous-$\mathcal{A}$-Modules de $\mathcal{A}$ seront qualifiés de
faisceaux d'idéaux (à gauche, à droite ou bilatères) dans $\mathcal{A}$ ou de
$\mathcal{A}$-Idéaux.

Lorsque $\mathcal{A}$ est un faisceau d'anneaux commutatifs, et que, dans la
définition des $\mathcal{A}$-Modules, on remplace partout la structure de
module par celle d'algèbre, on obtient la définition d'une
$\mathcal{A}$-Algèbre sur $X$. Il revient au même de dire qu'une
$\mathcal{A}$-Algèbre (non nécessairement commutative) est un
$\mathcal{A}$-Module $\mathcal{C}$, muni d'un homomorphisme de
$\mathcal{A}$-Modules
$\varphi:\mathcal{C}\otimes_{\mathcal{A}}\mathcal{C}\to\mathcal{C}$ et
d'une section $e$ au-dessus de $X$, tels que~: 1\textsuperscript{o} Le diagramme
